\begin{frame}[label=q5d4]{5d.4 · Funciones: construir fecha y trimestre}
  \small
  \pregunta{¿Cuántas personas se reportan por trimestre al construir una fecha con año y mes?}
  \entrada{93\,831 filas. \texttt{anio}, \texttt{mes}, \texttt{cantidad}. La fecha usa día 01 convencional.}
  \flujo{Año y mes originales}{\texttt{concat\_ws + lpad}\\\texttt{to\_date}}{\texttt{quarter(fecha)}\\Agrupar y sumar}{24 trimestres\\Orden cronológico}
  \begin{columns}[T,onlytextwidth]
    \column{0.55\textwidth}
    \textbf{Código Scala original (fragmento)}
    \codigo{sql/05d_4_fecha_trimestre}{12}{13}
    \codigo{sql/05d_4_fecha_trimestre}{18}{20}
    \column{0.42\textwidth}
    \textbf{Salida}\\[5pt]
    {\footnotesize
\begin{tabular}{@{}ll@{}}
\toprule
Trimestre & Personas\\
\midrule
2019-T4 & 7\,164\\
2020-T1 & 6\,018\\
2020-T2 & 2\,485\\
2020-T3 & 4\,567\\
\bottomrule
\end{tabular}
    }
  \end{columns}
  \respuesta{Entre los trimestres mostrados, 2020-T2 presenta el menor total. La consulta construye una clave temporal y describe la serie; no identifica por sí sola la causa de la caída.}
  \origen{sql}{05d_4_fecha_trimestre}
\end{frame}
